\begin{proof}[Proof of \cref{obj:lem:population-positivity}]
Fix \(v\in\mathbb R^q\). We compute the population entries, contract them with \(v\), and then obtain the energy decomposition and lower bound.

\begin{enumerate}
\item The localization density and fine-basis coordinates are
\[
w_h(x)=h^{-d}\mathbf 1_{C_h}(x),
\qquad
k=qJ^d,
\qquad
b_{J,z_{\mathrm f}}(x)=[\mathbf b_J(x)]_{z_{\mathrm f}},
\quad z_{\mathrm f}\in\{1,\ldots,k\}.
\]
Here \(x\in\mathbb R^d\), and the density is as in
\cref{obj:synth_3}. The role counts from \cref{obj:def:roles} satisfy
\[
\ell=\lfloor n/2\rfloor\ge1,
\qquad
t=n+m-\ell\ge n-\ell\ge1.
\]
For fixed \(h,J\), the localized basis products are bounded and measurable, so all expectations and integrals below are finite.

For coarse indices \(u,w\in\{1,\ldots,q\}\), expand the kernel in
\cref{obj:def:projection} inside the raw matrix of
\cref{obj:def:moments}. This gives
\[
\begin{aligned}
\widehat Q^{\mathrm{raw}}_{h,J;uw}
&=\frac1\ell\sum_{j\in I_L}
 w_h(X_j)r_u(X_j)r_w(X_j)A_j\\
&\quad-\sum_{z_{\mathrm f}=1}^{k}\frac1{t\ell}
 \sum_{i\in I_T}\sum_{j\in I_L}
 \bigl[w_h(X_i^T)r_u(X_i^T)b_{J,z_{\mathrm f}}(X_i^T)A_i^T\bigr]
 \bigl[w_h(X_j)r_w(X_j)b_{J,z_{\mathrm f}}(X_j)A_j\bigr].
\end{aligned}
\]
Under the product sampling law of \cref{obj:ass:sample-law}, each covariate-treatment pair has marginal law \(P_{XA}\). Each treatment-role record is independent of each outcome-role record: for a labeled treatment-role record their indices are distinct, and for an auxiliary record independence follows from the product sampling law. Taking expectations term by term therefore yields
\[
\begin{aligned}
\mathbb E_{\mathsf E_{n,m}(P)}
 [\widehat Q^{\mathrm{raw}}_{h,J;uw}]
&=\frac{\ell}{\ell}\,
 \mathbb E_P[w_h(X)r_u(X)r_w(X)A]\\
&\quad-\sum_{z_{\mathrm f}=1}^{k}\frac{t\ell}{t\ell}\,
 \mathbb E_P[w_h(X)r_u(X)b_{J,z_{\mathrm f}}(X)A]\,
 \mathbb E_P[w_h(X)r_w(X)b_{J,z_{\mathrm f}}(X)A].
\end{aligned}
\]
For any of the functions
\[
f_{\mathrm{loc}}\in
\{r_ur_w,\ r_ub_{J,z_{\mathrm f}},\ r_wb_{J,z_{\mathrm f}}\},
\qquad
f_{\mathrm{loc}}:\mathbb R^d\longrightarrow\mathbb R,
\]
the designated conditional treatment mean and uniform design in
\cref{obj:def:primitive-class} give
\[
\begin{aligned}
\mathbb E_P[w_h(X)f_{\mathrm{loc}}(X)A]
&=\mathbb E_P[w_h(X)f_{\mathrm{loc}}(X)e_P(X)]\\
&=h^{-d}\int_{C_h}f_{\mathrm{loc}}(x)e_P(x)\,dx\\
&=\int f_{\mathrm{loc}}(x)e_P(x)\,d\nu_h(x).
\end{aligned}
\]
The localization cube lies in \([0,1]^d\), so these identities use the designated witness on its specified domain. Substitution gives a formula symmetric in \(u,w\). Consequently, the symmetrization in \cref{obj:def:moments} preserves its expectation, and
\[
\begin{aligned}
(Q_{P,h,J})_{uw}
&=\int r_u(x)e_P(x)r_w(x)\,d\nu_h(x)\\
&\quad-\sum_{z_{\mathrm f}=1}^{k}
 \left(\int b_{J,z_{\mathrm f}}(x)e_P(x)r_u(x)\,d\nu_h(x)\right)
 \left(\int b_{J,z_{\mathrm f}}(x)e_P(x)r_w(x)\,d\nu_h(x)\right).
\end{aligned}
\]
% lean: CausalSmith.Stat.TwosamplePointcateAnnotationFrontier.population_gram_entries

\item Define the weighted coarse-basis coefficients by
\[
c_{z_{\mathrm f},u}
=\int b_{J,z_{\mathrm f}}(x)e_P(x)r_u(x)\,d\nu_h(x),
\qquad
1\le z_{\mathrm f}\le k,\quad 1\le u\le q.
\]
Multiplying the entry formula by \(v_uv_w\), summing, and using
\(p_v(x)=\sum_{u=1}^{q}r_u(x)v_u\), we obtain
\[
\begin{aligned}
v^\top Q_{P,h,J}v
&=\int e_P(x)
 \left(\sum_{u=1}^{q}r_u(x)v_u\right)^2\,d\nu_h(x)
 -\sum_{z_{\mathrm f}=1}^{k}
 \left(\sum_{u=1}^{q}c_{z_{\mathrm f},u}v_u\right)^2\\
&=\int e_P(x)p_v(x)^2\,d\nu_h(x)
 -\sum_{z_{\mathrm f}=1}^{k}
 \left(\sum_{u=1}^{q}c_{z_{\mathrm f},u}v_u\right)^2.
\end{aligned}
\]
Indeed, the integral term follows by expanding the finite double sum, and the correction follows from
\[
\sum_{u,w=1}^{q}v_uv_w
 \sum_{z_{\mathrm f}=1}^{k}c_{z_{\mathrm f},u}c_{z_{\mathrm f},w}
=
\sum_{z_{\mathrm f}=1}^{k}
 \left(\sum_{u=1}^{q}c_{z_{\mathrm f},u}v_u\right)^2.
\]

The finite kernel expansion in \cref{obj:def:projection} and linearity of integration give
\[
\begin{aligned}
\Pi_J(e_Pp_v)(x)
&=\sum_{z_{\mathrm f}=1}^{k}b_{J,z_{\mathrm f}}(x)
 \int b_{J,z_{\mathrm f}}(y)e_P(y)p_v(y)\,d\nu_h(y)\\
&=\sum_{z_{\mathrm f}=1}^{k}b_{J,z_{\mathrm f}}(x)
 \sum_{u=1}^{q}c_{z_{\mathrm f},u}v_u.
\end{aligned}
\]
The fine basis in that definition is orthonormal. Expanding the square and integrating therefore gives
\[
\begin{aligned}
\int \bigl[\Pi_J(e_Pp_v)(x)\bigr]^2\,d\nu_h(x)
&=\sum_{z_{\mathrm f},z'_{\mathrm f}=1}^{k}
 \left(\sum_{u=1}^{q}c_{z_{\mathrm f},u}v_u\right)
 \left(\sum_{w=1}^{q}c_{z'_{\mathrm f},w}v_w\right)
 \int b_{J,z_{\mathrm f}}b_{J,z'_{\mathrm f}}\,d\nu_h\\
&=\sum_{z_{\mathrm f}=1}^{k}
 \left(\sum_{u=1}^{q}c_{z_{\mathrm f},u}v_u\right)^2.
\end{aligned}
\]
Thus
\[
v^\top Q_{P,h,J}v
=\int e_Pp_v^2\,d\nu_h
 -\int \bigl[\Pi_J(e_Pp_v)\bigr]^2\,d\nu_h.
\]
% lean: CausalSmith.Stat.TwosamplePointcateAnnotationFrontier.population_gram_quadratic_of_entries

\item The finite coarse expansion makes \(p_v\) square-integrable, and bounded overlap makes \(e_Pp_v\) square-integrable. The projection expansion just established also yields
\[
\begin{aligned}
\int e_Pp_v\,\Pi_J(e_Pp_v)\,d\nu_h
&=\sum_{z_{\mathrm f}=1}^{k}
 \left(\int b_{J,z_{\mathrm f}}e_Pp_v\,d\nu_h\right)^2\\
&=\int \bigl[\Pi_J(e_Pp_v)\bigr]^2\,d\nu_h.
\end{aligned}
\]
Expanding the residual square and using this equality gives the Pythagorean identity
\[
\int (e_Pp_v)^2\,d\nu_h
=
\int \bigl[\Pi_J(e_Pp_v)\bigr]^2\,d\nu_h
+\int \bigl[e_Pp_v-\Pi_J(e_Pp_v)\bigr]^2\,d\nu_h.
\]
Meanwhile, the pointwise identity
\(e_P(1-e_P)p_v^2=e_Pp_v^2-(e_Pp_v)^2\) gives
\[
\int e_P(1-e_P)p_v^2\,d\nu_h
=
\int e_Pp_v^2\,d\nu_h-\int(e_Pp_v)^2\,d\nu_h.
\]
Combining these two identities with the quadratic-form formula proves
\[
v^\top Q_{P,h,J}v
=
\int e_P(1-e_P)p_v^2\,d\nu_h
+\int\bigl[e_Pp_v-\Pi_J(e_Pp_v)\bigr]^2\,d\nu_h.
\]
% lean: CausalSmith.Stat.TwosamplePointcateAnnotationFrontier.population_projection_energy_identity

\item The scaled coarse polynomials of \cref{obj:synth_1} satisfy
\[
\int r_u(x)r_w(x)\,d\nu_h(x)
=
\begin{cases}
1,&u=w,\\
0,&u\ne w.
\end{cases}
\]
To see this, integrate their tensor products over the localization cube and rescale each coordinate to \([-1/2,1/2]\). Direct integration of the three normalized one-dimensional polynomials in that definition gives unit squared norms and zero pairwise inner products. When \(u=w\), every coordinate factor is one; when \(u\ne w\), some multi-index coordinate differs and its factor is zero. Hence finite expansion gives
\[
\int p_v(x)^2\,d\nu_h(x)
=\sum_{u,w=1}^{q}v_uv_w\int r_ur_w\,d\nu_h
=\sum_{u=1}^{q}v_u^2.
\]

Class membership supplies \cref{obj:ass:overlap}. Since \(\nu_h\) is supported on \(C_h\subseteq[0,1]^d\), for \(\nu_h\)-almost every \(x\),
\[
e_P(x)(1-e_P(x))-\varepsilon(1-\varepsilon)
=
\left(e_P(x)-\varepsilon\right)
\left(1-\varepsilon-e_P(x)\right)
\ge0.
\]
Multiplying by \(p_v(x)^2\) and integrating therefore yields
\[
\int e_P(1-e_P)p_v^2\,d\nu_h
\ge\varepsilon(1-\varepsilon)\int p_v^2\,d\nu_h
=\varepsilon(1-\varepsilon)\sum_{u=1}^{q}v_u^2.
\]
Finally,
\[
\int\bigl[e_Pp_v-\Pi_J(e_Pp_v)\bigr]^2\,d\nu_h\ge0,
\]
so the decomposition from the preceding step implies
\[
v^\top Q_{P,h,J}v\ge\varepsilon(1-\varepsilon)\sum_{u=1}^{q}v_u^2.
\]
% lean: CausalSmith.Stat.TwosamplePointcateAnnotationFrontier.population_expression_positivity
\end{enumerate}
\end{proof}
